\section{Introduction}
An operating contract is evaluated through the decisions it induces. A guarantee can discourage surrender while reducing the agent's private value; a compulsory operating term can substitute for a guarantee while imposing a real enforcement cost. When consumption, financial exposure, deliberate preference adjustment, and surrender are chosen jointly, comparing private values does not establish a comparison of purchased service. Nor does approximating one optimal policy establish the service delivered by every optimal policy. At an indifference point an arbitrarily small value error can conceal an economically important change in the response.

This paper develops Neural Bellman Operators as a proposal, evaluation, improvement, and verification architecture for this problem. The economic object passed from a dynamic program to the purchaser is a joint response set, not a selected policy or an unqualified numerical derivative. Certified reward perturbations constrain service, surrender, and quality simultaneously. A purchaser compares the least favorable response to an executable offer with upper bounds on competing offers, after charging participation transfers and enforcement costs. Neural proposals and their independent certificates perform distinct tasks: a proposal searches for a useful feasible policy; the certificate determines what can be concluded about all responses and rival contracts.

Our first economic result concerns the information required by that comparison. Theorem~\ref{thm:r19-information} constructs two finite dynamic economies that agree on an entire family of value queries, including every positive and negative coordinate query, but imply opposite procurement decisions. More coordinate accuracy, and even adaptive acquisition of infinitely many such messages, cannot resolve the ranking. One query in the purchaser's joint direction can. The theorem gives its exact adversarial-noise threshold, the minimax regret without that direction, and the expected value of the information under any prior on the two economies. The institution enters through the outputs accruing to the purchaser: when one of those outputs ceases to matter, the same information can become worthless. Proposition~\ref{prop:r19-sharp} complements this separation with a sharp general response set under exact messages. These are information results for a value-oracle purchaser, not lower bounds against an algorithm already supplied with the complete transition model.

The paper then separates implementation at an initial contract from implementation at every state. A larger operating action set raises the no-surrender value state by state and weakly lowers a uniform stopping-obstacle requirement. It need not lower the initial requirement without additional conditions. Proposition~\ref{prop:r19-initial-order} supplies two such conditions: reachability under a no-surrender optimum, or a response-preserving comparison of the adjustment gain on surrender deviations and on continued operation. Theorem~\ref{thm:r19-enforcement} identifies the initial fee threshold and gives a sharp global near-optimal-response frontier. A fee slack $F-F^{\rm init}$ bounds surrender by $\eta/(F-F^{\rm init})$. This distinction matters when a purchaser wants a guarantee for every approximately optimal response rather than for a selected exact policy.

A third result connects continuous instruments to that global behavior. Endpoint policies generate lower affine payoff planes, and endpoint action values generate upper chords over a whole charge interval. Their verified separation supplies a lower bound on each action's Bellman loss. Theorem~\ref{thm:r19-cell} combines these bounds in one occupancy-weighted loss budget. At zero tolerance it gives a whole-interval exclusion graph; at positive tolerance it retains rare, locally costly actions when their total expected loss is small. The construction is explicitly an outer relaxation: actions that require different charges may survive it, but their coexistence is never presented as an implementable common-charge response.

The continuous-state and quantitative arguments are kept distinct. Proposition~\ref{prop:r19-continuous} exhibits a controlled Brownian economy with explicit continuation values, an admissible feedback, and globally verified stopping obstacles. Its initial enforcement thresholds are $(\log2)/2$ with adjustment and $1/2$ without adjustment. Rational enclosures establish a strict procurement comparison, including participation and positive-tolerance response error. This is an executed continuous verification chain, not a small sampled residual. The original stochastic settlement application retains its CRRA operating primitive, terminal settlement, financial mandates, and all adverse cases. Its frozen arrays are a separate numerical target. The analytic example does not supply the missing constructor bound for that target's antecedent diffusion.

The settlement evidence shows why institutional precision matters. On the original enforcement menu the preferred compulsory terms are three intervals with adjustment and four without; refinement to increments of 0.05 gives two versus four; continuous charges identify term one in both regimes. These are different feasible contract sets, not conflicting estimates of the same optimization problem. The complete continuous-fee cover, rational upper certificates, and executable lower offers are retained. New threshold calculations on the same arrays independently bracket the initial no-surrender fees under both initial financial mandates. Lower endpoints have feasible profitable surrender deviations; upper endpoints have a verified all-optimal-action graph from which surrender is unreachable. At the common charge 0.765, adjustment implements zero elective surrender under exact responses, whereas the no-adjustment contract retains a strictly positive surrender option. None of these claims relies on the uniform-threshold ordering alone.

The computational evidence also retains its unfavorable outcomes. The original network, polynomial policy, nearest-anchor policy, common-law upper methods, and parameter-lifting comparison remain visible. A new three-seed experiment compares the trained network with nearest-anchor, polynomial, and fully evaluated teacher-bank controls on the unchanged target; every method is charged its common exact upper reference and its policy evaluation. Fitting, teachers, online work, and independent checking are reported separately. This experiment studies repeatability and the cost of certified decisions. It does not vary state dimension, and it does not establish an end-to-end neural advantage over exact dynamic programming. The operator theory and information results are not conditioned on winning that benchmark.

\paragraph*{Relation to established methods.}
Occupancy measures, performance-difference identities, Lagrangian support duality, and Bellman verification are established mathematical tools. We use them as tools rather than presenting their individual identities as new economic theorems. The contribution isolated here is the institutional information requirement and the implementation conclusions obtained after global response uncertainty is retained. The query experiment proves a strict information separation and its decision cost; the implementation theorems identify when a uniform ordering reaches the initial contract and how fee slack prices behavioral error. The exact diffusion example supplies the candidate and decision budget required by verification, while the finite settlement economy tests the institutional comparisons without deleting their adverse cases.

\paragraph*{Reading the argument.}
The next section specifies the economic environment. The information and procurement sections distinguish message-based identification from known-model restrictions. The enforcement and interval sections develop implementation guarantees for exact and globally near-optimal agents. The continuous benchmark and settlement results then give separately identified applications. Full proofs, intermediate derivations, arithmetic details, and additional computational evidence are in the supplementary appendix. Earlier operator, recursive-utility, equilibrium, and mechanism derivations remain in the unchanged historical compendium; the preservation map identifies their location without inserting repository chronology into the argument.
